% The class cannot take \input inside \abstract{}, so the text is a macro.
\newcommand{\AbstractText}{%
Trained seizure-alert assistance dogs signal an oncoming seizure, but the signal helps only if a human perceives it---and the owner is the person least able to. A recent study showed that a collar-mounted inertial sensor can detect a trained spin alert; it stopped at detection. We build and evaluate the rest: an ESP32 collar with a streaming $K$-of-$N$ verification gate, a phone relay running an owner $\rightarrow$ caregiver $\rightarrow$ emergency escalation ladder, and a backend watchdog. On the public dataset (\R{ds.spins} alerts, six dogs) the deployed forest reaches leave-one-group-out event recall \R{logo.rfe.evRec} (95\% CI \R{logo.rfe.evRecCI}). Scored on continuous recordings with nested operating-point selection, however, it raises \R{cont.rfe.r0.fah} false alarms per hour, \R{cont.ratio.curated}$\times$ the curated-clip estimate. We specified the escalation protocol in TLA$^+$: \R{tlc.nchecks} model-checked properties confirmed its safety, exposed a watchdog flaw that bypasses the caregiver when the phone dies or loses connectivity (fixed and re-verified), and led to an integration bug that disabled the watchdog. A Monte-Carlo simulation driving the real engine shows that escalation reaches a human within five minutes for \R{sim.ownerincapacitated.p300givendet}\% of detected alerts, and that limiting caregiver pages to one per week requires at most \R{req.maxfah} false alarms per hour---\R{req.factor}$\times$ below the measured rate. Escalation can be verified; detection is the bottleneck.
}
